\documentclass[reqno]{amsart}
\usepackage{amsmath,amssymb,amsthm,mathtools}
\usepackage[margin=1in]{geometry}
\usepackage{enumitem}
\setlist{nosep,leftmargin=*}
\usepackage[colorlinks=true,linkcolor=blue,citecolor=blue]{hyperref}

\theoremstyle{plain}
\newtheorem{theorem}{Theorem}[section]
\newtheorem{proposition}[theorem]{Proposition}
\newtheorem{lemma}[theorem]{Lemma}
\newtheorem{corollary}[theorem]{Corollary}
\theoremstyle{definition}
\newtheorem{definition}[theorem]{Definition}
\newtheorem{remark}[theorem]{Remark}
\newtheorem{question}[theorem]{Question}

\DeclareMathOperator{\rpp}{rpp}
\newcommand{\dk}{d_k}

\begin{document}

\title{The skew hook formula of Naruse--Okada does not extend to
simple-spectrum heaps}

\author{Jaideep Sai Padhi}
\address{Department of Computer Science, Purdue University}
\email{jpadhi@purdue.edu}

\date{\today}

\begin{abstract}
Let $P$ be the coloured poset attached to a Weyl group element $w$ for which
the torus acts with simple spectrum, in the sense of
\cite[\S5.4.2]{DKK}.  Every such $P$ is a hook length poset, so its
reverse-plane-partition series is a hook product; it need not be $d$-complete.
We answer, negatively, the question of whether the skew hook formula of
Naruse--Okada \cite[Thm.~1.2]{NO} extends to this larger class.  The
obstruction is not the sum over excited diagrams but the hook function itself,
and it is already visible at $F=\varnothing$, where the formula degenerates to
the Peterson--Proctor product.  We exhibit an eight-element witness in $D_4$:
its true hook multiset is $\{1,1,3,4,4,5,6,7\}$, and the recursion of
\cite[Def.~2.6]{NO} --- which off the $d$-complete locus is a \emph{family} of
recursions, one per choice of $\dk$-interval --- produces none of
$\{1,1,3,4,4,5,6,7\}$ under any of its choices, in either orientation
convention.  The verification is finite and small enough to check by hand:
eight elements, at most eight choices.  We emphasise what is \emph{not}
disproved: since $P$ is a hook length poset, the $F=\varnothing$ case of the
formula is true for $P$ with the correct hook function, so the question of a
corrected recursion remains open, and we state it.
\end{abstract}

\maketitle

\section{The question, and the answer}

Fix a simply-laced Dynkin diagram $\Gamma$ with vertex set $I$ and Weyl group
$W$, a dominant weight $\lambda$, and $w\in W$.  Following
\cite[\S5.1]{DKK}, a reduced word for $w$ has a \emph{heap}: the poset on
$\{1,\dots,\ell\}$ generated by the relations between positions carrying equal
or adjacent colours, with elements coloured by $\pi(a)=i_a$.  When $w$ is
$\lambda$-minuscule the heap is $d$-complete, and
\cite[Thm.~1.1]{NO} --- the Peterson--Proctor hook-product formula --- computes
its reverse-plane-partition series.  Dinkins, Karpov and Krylov observe
\cite[\S5.4.2]{DKK} that the relevant hypothesis is weaker than
$\lambda$-minusculity: what is needed is that the torus act on each $V_i$ with
\emph{simple spectrum}, and they ask for a combinatorial characterisation of
such $w$.  For every such $w$ the resulting coloured poset $P=P(w)$ is conjectured, and
in a companion manuscript proved, to satisfy
\begin{equation}\label{eq:hook}
   \sum_{\varphi\in\rpp(P)}q^{|\varphi|}
   \;=\;\prod_{u\in P}\frac{1}{1-q^{\,h(u)}}
\end{equation}
for a function $h:P\to\mathbb{Z}_{>0}$ --- that is, $P$ is a \emph{hook length
poset} in the sense of Ishikawa--Tagawa \cite{IT} --- and the class of such $P$
strictly contains the $d$-complete ones.  Nothing below depends on that
general statement: the present note needs \eqref{eq:hook} only for the single
poset of \S3, where it is verified directly in Lemma \ref{lem:true}.

The question addressed here was put to the author by V.~Krylov in October
2026.

\begin{question}[Krylov]\label{q:main}
Naruse and Okada prove a \emph{skew} hook formula \cite[Thm.~1.2]{NO}: for $P$
$d$-complete and $F\subseteq P$ an order ideal, the series of reverse plane
partitions on $P\setminus F$ is a finite sum over excited diagrams of $F$ in
$P$, with hook function given by \cite[Def.~2.6]{NO}.  Does this formula extend
to the skew-heaps $P\setminus F$ arising in the simple-spectrum setting, where
$P$ is a hook length poset but need not be $d$-complete?
\end{question}

\begin{theorem}\label{thm:main}
No.  There is a heap $P$ with eight elements, arising from a reduced word in
$D_4$, which is a hook length poset, is not $d$-complete, and for which
\eqref{eq:hook} fails for \emph{every} hook function produced by
\cite[Def.~2.6]{NO}.  Consequently the skew formula \cite[Thm.~1.2]{NO}, read
with its own hook function, is false for $P$ already at $F=\varnothing$.
\end{theorem}

The witness is Proposition \ref{prop:witness}; the two finite verifications
are Lemma \ref{lem:true} and Proposition \ref{prop:exhaust}.  Section
\ref{sec:scope} says precisely what the theorem does and does not rule out,
and \S\ref{sec:open} states the question that replaces Question \ref{q:main}.

\section{The reduction to \texorpdfstring{$F=\varnothing$}{F empty}}

We use only the degenerate case of the skew formula.  As
\cite[\S5.3]{DKK} note in quoting it: when $F=\varnothing$, so that
$P\setminus F=P$, the set of excited diagrams consists of the empty diagram
alone, the set of excited peaks is empty, and the formula reduces to
\eqref{eq:hook}.  So if \cite[Thm.~1.2]{NO} holds for $P$ with hook function
$h$, then \eqref{eq:hook} holds with that $h$.

It therefore suffices to show that no $h$ permitted by \cite[Def.~2.6]{NO}
satisfies \eqref{eq:hook}.  The next lemma is what makes this a statement about
\emph{multisets}, and hence finitely checkable.

\begin{lemma}\label{lem:unique}
Let $(a_i)_{i\in A}$ and $(b_j)_{j\in B}$ be finite multisets of positive
integers with
$\prod_{i}(1-q^{a_i})^{-1}=\prod_{j}(1-q^{b_j})^{-1}$ in
$\mathbb{Z}[[q]]$.  Then the multisets coincide.
\end{lemma}

\begin{proof}
Write $c_d$ for the multiplicity of $d$ in $(a_i)$.  Taking logarithms,
\[
  \log\prod_i\frac{1}{1-q^{a_i}}
  =\sum_{m\ge1}\frac1m\sum_i q^{a_im}
  =\sum_{N\ge1}q^N\sum_{d\mid N}c_d\,\frac{d}{N},
\]
so $N\,[q^N]\log(\cdot)=\sum_{d\mid N}d\,c_d$ for every $N\ge1$.  The series
therefore determines $\sum_{d\mid N}d\,c_d$ for all $N$, and these determine
$d\,c_d$, hence $c_d$, by induction on $N$ (equivalently by M\"obius
inversion).  The same holds for $(b_j)$, and the two sets of data agree.
\end{proof}

Accordingly we call the multiset $\{h(u):u\in P\}$ of any $h$ satisfying
\eqref{eq:hook} the \emph{true hook multiset} of $P$; by Lemma
\ref{lem:unique} it is well defined, and it is an invariant of the series
alone.

\section{The witness}

\begin{proposition}\label{prop:witness}
Let $\Gamma=D_4$ in the Bourbaki labelling, so that vertex $2$ is trivalent
with neighbours $1,3,4$, and let
\[
   w \;=\; s_2\,s_1\,s_2\,s_3\,s_4\,s_2\,s_1\,s_3 ,
\]
a reduced word of length $8$.  Its heap $P$ has eight elements
$p_1\prec\cdots$ with colours and order
\[
  \begin{array}{c|cccccccc}
    \text{position} & p_1&p_2&p_3&p_4&p_5&p_6&p_7&p_8\\\hline
    \text{colour}   & 2 & 1 & 2 & 3 & 4 & 2 & 1 & 3
  \end{array}
  \qquad
  p_1\prec p_2\prec p_3\prec\{p_4,p_5\}\prec p_6\prec\{p_7,p_8\},
\]
with $p_4,p_5$ incomparable \textup{(}colours $3,4$ are non-adjacent\textup{)}
and $p_7,p_8$ incomparable \textup{(}colours $1,3$ are
non-adjacent\textup{)}.  Then $P$ is not $d$-complete.
\end{proposition}

\begin{proof}
The order is the transitive closure of the heap relations, read off from the
word: consecutive equal or adjacent colours are comparable, and the two
displayed pairs carry non-adjacent colours.  For the failure of
$d$-completeness, no computation is needed: $p_7$ and $p_8$ are incomparable
and have the common lower cover $p_6$, so they span a $d_3^-$-interval, and
Proctor's axiom at $k=3$ requires such an interval to have a top, i.e.\ a
common upper cover.  Both $p_7$ and $p_8$ are maximal in $P$, so they have no
upper bound at all.  (The dual poset is not $d$-complete either, so the
conclusion is independent of the orientation convention; see Remark
\ref{rem:orient}.)
\end{proof}

\begin{lemma}\label{lem:true}
The true hook multiset of the poset $P$ of Proposition \ref{prop:witness} is
\[
   \{1,1,3,4,4,5,6,7\}.
\]
\end{lemma}

\begin{proof}
Both sides of \eqref{eq:hook} are computed independently and compared to order
$q^{26}$.  For the left-hand side, a reverse plane partition
$\varphi:P\to\mathbb{Z}_{\ge0}$ with $\varphi(x)\le\varphi(y)$ for
$x\prec y$ is equivalent to the decreasing chain of order filters
$F_k=\{x:\varphi(x)\ge k\}$, with $|\varphi|=\sum_{k\ge1}|F_k|$, so
\[
  \sum_{\varphi\in\rpp(P)}q^{|\varphi|}
  \;=\;1+\sum_{\varnothing\ne F}g(F),
  \qquad
  g(F)=\frac{q^{|F|}}{1-q^{|F|}}\Bigl(1+\sum_{\varnothing\ne G\subsetneq F}g(G)\Bigr),
\]
the second identity because repetitions $F_k=F_{k+1}$ are permitted.  $P$ has
ten nonempty order filters, so this is a ten-term recursion.  It gives
\[
  1+2q+3q^2+5q^3+9q^4+14q^5+21q^6+31q^7+45q^8+63q^9+87q^{10}+118q^{11}+\cdots,
\]
which agrees with $\prod_{h\in\{1,1,3,4,4,5,6,7\}}(1-q^h)^{-1}$ in every
coefficient up to $q^{26}$.  By Lemma \ref{lem:unique} the multiset is the
only one that can occur.  As a sanity check on the first two coefficients by
hand: $[q^1]=2$ counts the two maximal elements $p_7,p_8$, and $[q^2]=3$
counts $\varphi=2\cdot\mathbf 1_{p_7}$, $2\cdot\mathbf 1_{p_8}$ and
$\mathbf 1_{p_7}+\mathbf 1_{p_8}$.
\end{proof}

\section{Exhausting Definition 2.6}

Off the $d$-complete locus an element of $P$ may be the top of more than one
$\dk$-interval, so \cite[Def.~2.6]{NO} does not define a single function: it
defines one for each choice of $\dk$-interval at each such element, together
with the option of using the order-ideal count $\#\{v:v\preceq u\}$ there
instead.  We exhaust that family.

\begin{proposition}\label{prop:exhaust}
For $P$ as in Proposition \ref{prop:witness}:
\begin{enumerate}
\item in the orientation of Proposition \ref{prop:witness}, three elements top
a $\dk$-interval, giving eight choices and six distinct multisets,
\[
  \{1,2,3,4,4,5,6,6\},\ \{1,2,3,4,4,5,6,7\},\ \{1,2,3,4,4,5,7,7\},
\]
\[
  \{1,2,3,4,4,6,6,6\},\ \{1,2,3,4,4,6,6,7\},\ \{1,2,3,4,4,6,7,7\};
\]
\item in the opposite orientation --- the corner-as-maximum convention in which
\cite[Def.~2.6]{NO} is stated --- two elements top a $\dk$-interval, giving six
choices and two distinct multisets,
\[
  \{1,1,3,4,4,5,7,8\},\qquad \{1,1,3,4,4,6,7,8\}.
\]
\end{enumerate}
In neither case does $\{1,1,3,4,4,5,6,7\}$ occur.
\end{proposition}

\begin{proof}
Finite enumeration; eight and six cases respectively.
\end{proof}

\begin{proof}[Proof of Theorem \ref{thm:main}]
Suppose \cite[Thm.~1.2]{NO} held for $P$ with hook function $h$ given by
\cite[Def.~2.6]{NO}.  Taking $F=\varnothing$ gives \eqref{eq:hook} for $h$, so
by Lemma \ref{lem:unique} the multiset of $h$ is the true hook multiset of
$P$, which by Lemma \ref{lem:true} is $\{1,1,3,4,4,5,6,7\}$.  Proposition
\ref{prop:exhaust} shows no choice in the family defining $h$ has that
multiset, in either orientation.
\end{proof}

\begin{remark}\label{rem:orient}
Both orientations are recorded because the hook function is pinned where the
series is a hook product, while \cite[Def.~2.6]{NO} is stated in the dual
convention, and the correspondence between the two is exactly the kind of
bookkeeping in which an error would hide.  The conclusion does not depend on
the choice.
\end{remark}

\begin{remark}
The failure in the orientation of (i) is close: $\{1,2,3,4,4,5,6,7\}$ differs
from the true multiset in one entry, $2$ against $1$.  In the orientation of
(ii) it is gross: the two largest hooks come out $7,8$ instead of $6,7$, so
the product has the wrong behaviour at $q\to1$ and no cancellation can repair
it.
\end{remark}

\section{What is, and is not, disproved}\label{sec:scope}

It is worth being exact, because the negative statement is narrower than
``the skew formula fails'' and the residue is the interesting part.

\begin{enumerate}
\item \emph{What fails} is \cite[Thm.~1.2]{NO} read with its own hook
function, on hook length heaps that are not $d$-complete.  The obstruction is
the hook function, and it is present at $F=\varnothing$, before excited
diagrams play any role.
\item \emph{What does not fail} is \eqref{eq:hook} itself.  The poset $P$ of
Proposition \ref{prop:witness} \emph{is} a hook length poset, by Lemma
\ref{lem:true}; the $F=\varnothing$ case of the formula is true for $P$, with
the hook function realising $\{1,1,3,4,4,5,6,7\}$.  So nothing here suggests
that the shape of \cite[Thm.~1.2]{NO} is wrong for skew-heaps.  What is wrong
is the recipe for $h$.
\item \emph{Nothing is claimed about} skew $F\ne\varnothing$ beyond what (i)
forces, nor about any formula whose hook function is defined otherwise, nor
about $d$-complete $P$, where \cite[Thm.~1.2]{NO} is a theorem and is not in
question.
\end{enumerate}

For context on how typical the witness is: over all reduced words of length at
most $8$ in $D_4$, among the heaps that are hook length posets, the
$d$-complete ones reach their true hook multiset in all $336$ cases, and the
non-$d$-complete ones fail to reach it in all $36$.  So the dichotomy in this
range is exactly $d$-completeness.  All $36$ failures have $|P|=8$; the
witness of Proposition \ref{prop:witness} is one of them, and $8$ is the
smallest size at which a failure occurs in this range.

\section{The question that replaces it}\label{sec:open}

\begin{question}\label{q:new}
Every simple-spectrum heap $P$ is a hook length poset, so it carries a hook
function $h_P$ satisfying \eqref{eq:hook}, unique as a multiset by Lemma
\ref{lem:unique}.  Is there a recursion computing $h_P$ --- agreeing with
\cite[Def.~2.6]{NO} on the $d$-complete locus --- for which the skew formula
\cite[Thm.~1.2]{NO} holds on all of this class?
\end{question}

Theorem \ref{thm:main} says that the recursion of \cite[Def.~2.6]{NO} is not
such a recursion.  It leaves open that one exists, and item (ii) of
\S\ref{sec:scope} is weak evidence that it might: the $F=\varnothing$ case,
which is the only case we have tested against, holds throughout the class once
$h$ is chosen correctly.

\section{Reproduction}

All claims above are finite checks.  The computations are available with the
rest of the project at
\url{https://github.com/jaideepsaipadhi/dkk-simple-spectrum}; in the
\texttt{skewheap} directory,
\begin{itemize}
\item \verb|python3 certify2.py| computes the left-hand side of
\eqref{eq:hook} for the witness by the filter recursion of Lemma
\ref{lem:true}, in both orientations, and compares with the hook product ---
from scratch, sharing no code with the rest of the project;
\item \verb|python3 orient.py| exhausts the \cite[Def.~2.6]{NO} family in both
orientations, printing every reachable multiset (Proposition
\ref{prop:exhaust});
\item \verb|python3 minwit.py 8 D4| searches for the smallest witness and
recovers the one above;
\item \verb|python3 reach.py 8 11 D4| produces the $336$ and $36$ of
\S\ref{sec:scope}.
\end{itemize}

\subsection*{Acknowledgements}
The question answered here was posed to the author by Vasily Krylov, who also
formulated the surrounding problem and discussed its details.  The
hook-product identity for simple-spectrum heaps, of which \eqref{eq:hook} is
the statement used here, is a conjecture of Dinkins, Karpov and Krylov.

\begin{thebibliography}{99}

\bibitem{DKK} H.~Dinkins, I.~Karpov and V.~Krylov, \emph{The quantum Hikita
conjecture via quasimaps}, arXiv:2608.16746.

\bibitem{IT} M.~Ishikawa and H.~Tagawa, \emph{Leaf posets and multivariate
hook length property}, RIMS K\^oky\^uroku \textbf{1913} (2014), 67--80.

\bibitem{NO} H.~Naruse and S.~Okada, \emph{Skew hook formula for $d$-complete
posets via equivariant $K$-theory}, Algebr.\ Comb.\ \textbf{2} (2019),
no.~4, 541--571.

\end{thebibliography}

\end{document}
